\section{Privacy, reliability, and quality assurance}

\subsection{Permission and privacy model}
Questory requests only permissions connected to a visible action. Location is
explained immediately before the system prompt; camera use begins from a quest
action. Precise tracks, captions, and retained image paths are not intentionally
logged. No account is required. Online sharing is opt-in, visible, time-limited,
and restricted to a summarized payload rather than the editable local database.
An explicit confirmation describes the precise route and media that will leave
the device before any network request begins.

The app never promises that a public route is safe. Bundled notes remind the user
to check current pedestrian conditions, weather, traffic, lighting, and local
rules. Pause, finish, discard, and stop-sharing controls must remain discoverable.

\subsection{Failure handling}
The implementation treats expected failures as product states: unavailable city
content, denied permission, disabled location service, camera cancellation,
missing retained media, storage failure, export failure, remote failure, and
empty History. Network failure cannot cancel or alter a local run. Zero-distance
and missing-photo cases have safe representations instead of crashes.

\subsection{Automated verification}
\begin{table}[htbp]
\centering
\small
\begin{tabularx}{\textwidth}{@{}p{40mm}p{29mm}X@{}}
\toprule
\textbf{Verification} & \textbf{Result} & \textbf{Meaning} \\
\midrule
\code{flutter analyze} & \status{Passed} & No static-analysis issues in the
verified repository state \\
Complete Flutter test suite & \status{58 passed} & Domain, controller,
persistence, widget, navigation, and export behaviors passed \\
Debug web compilation & \status{Passed} & Shared Flutter code compiles; this
does not claim production browser persistence support \\
LaTeX build and log scan & \status{Passed} & Report compiles reproducibly; final
page render is visually inspected after screenshots are inserted \\
Android minimum SDK & \status{API 24} & The application build explicitly
supports the required minimum Android level \\
\bottomrule
\end{tabularx}
\caption{Verified automated evidence as of 5 September 2026.}
\end{table}


Tests cover distance calculations and invalid-jump filtering, active versus
paused duration, zero-distance pace, quest proximity and fallback, achievement
evaluation, persistence round-trips, story transforms and serialization,
undo/redo, responsive screens, permission recovery, navigation, and export
feedback. Platform plugins are wrapped so these behaviors are exercised with
deterministic repositories, trackers, clocks, photo stores, renderers, and share
services.

\subsection{Visual and accessibility review}
Core screens were rendered at a representative Android portrait size for the
figures in this report. The review checks clipped text, overflow, contrast,
control reachability, status meaning, empty/error copy, and separation between
editor controls and exported story content. Real-device review remains important
for system permission dialogs, camera return behavior, notification appearance,
and continuous GPS under background conditions.

\subsection{Engineering limitations}
The product deliberately excludes turn-by-turn navigation, public rankings,
weather, accounts, cloud synchronization, and Canva handoff. These would expand
permissions, availability risks, and privacy obligations without improving the
offline academic flow. The codebase also contains a focused Story Studio browser
harness; it is development support and must not be confused with the Android
production entry point.
